% id: 54-e6
% book: 개념원리 확률과 통계 · 05 이항정리의 활용
% section: 필수·발전 예제
% figure: none
% answer: ⑴ $7$ ⑵ $1$ ⑶ $4$ ⑷ $1023$
% answer_source: 본문 풀이
% variant_level: 0 (원본 전사)
% 이 파일은 build.mjs 가 items.json 에서 만든다. 고칠 때는 items.json 을 고친다.
\smash{\raisebox{1.5mm}{\dmkichul{개념원리 확률과 통계 54쪽 예제 6 · 필수}}}
다음을 구하시오.
\dmsubprob{1} ${}_7\mathrm{C}_0+{}_7\mathrm{C}_1+{}_7\mathrm{C}_2+\cdots+{}_7\mathrm{C}_7=2^n$을 만족시키는 자연수 $n$의 값
\dmsubprob{2} ${}_{20}\mathrm{C}_1-{}_{20}\mathrm{C}_2+{}_{20}\mathrm{C}_3-{}_{20}\mathrm{C}_4+\cdots+{}_{20}\mathrm{C}_{19}-{}_{20}\mathrm{C}_{20}$의 값
\dmsubprob{3} ${}_{2n}\mathrm{C}_0+{}_{2n}\mathrm{C}_2+{}_{2n}\mathrm{C}_4+\cdots+{}_{2n}\mathrm{C}_{2n}=128$을 만족시키는 자연수 $n$의 값
\dmsubprob{4} ${}_{11}\mathrm{C}_1+{}_{11}\mathrm{C}_3+{}_{11}\mathrm{C}_5+{}_{11}\mathrm{C}_7+{}_{11}\mathrm{C}_9$의 값
